\chapter*{LỜI MỞ ĐẦU}
\addcontentsline{toc}{chapter}{Lời mở đầu}

\section*{Lý do lựa chọn đề tài}

Tín dụng là hoạt động cốt lõi của ngân hàng thương mại, đồng thời là nguồn phát sinh rủi ro có ảnh hưởng trực tiếp đến lợi nhuận, chất lượng tài sản và mức độ an toàn của ngân hàng. Vì vậy, đánh giá và xếp hạng tín dụng không chỉ phục vụ quyết định cấp tín dụng mà còn hỗ trợ giám sát khoản vay, phân loại nợ và nhận diện sớm dấu hiệu suy giảm chất lượng tín dụng.

Sự gia tăng về quy mô và mức độ đa dạng của dữ liệu ngân hàng đặt ra yêu cầu đổi mới phương pháp đánh giá rủi ro. Các phương pháp truyền thống có ưu điểm về tính minh bạch, nhưng có thể gặp hạn chế trước quan hệ phi tuyến, tương tác phức tạp và tình trạng mất cân bằng giữa các nhóm nợ. Trí tuệ nhân tạo và học máy vì vậy mở ra khả năng khai thác dữ liệu hiệu quả hơn, đồng thời hỗ trợ chuẩn hóa và tự động hóa một phần quy trình đánh giá.

Trong bối cảnh đó, đề tài \textit{``Ứng dụng trí tuệ nhân tạo trong xếp hạng tín dụng: nghiên cứu từ một ngân hàng thương mại ở Việt Nam''} được thực hiện nhằm đánh giá khả năng áp dụng các mô hình học máy vào phân loại năm nhóm nợ trên dữ liệu thực tế. Nghiên cứu đồng thời xem xét vai trò của các nhóm thông tin nghiệp vụ và khả năng rút gọn đặc trưng, qua đó tạo cơ sở cho việc phát triển hệ thống cảnh báo sớm rủi ro tín dụng.

Phần thực nghiệm hiện tại tập trung vào phân loại nhóm nợ và đánh giá mô hình ngoài mẫu. Dự báo khả năng chuyển nhóm nợ theo thời gian được xác định là hướng phát triển tiếp theo khi dữ liệu có đầy đủ lịch sử trạng thái và cửa sổ dự báo.

\section*{Đối tượng và phạm vi nghiên cứu}

Đối tượng nghiên cứu là việc ứng dụng các phương pháp thống kê và học máy trong phân loại mức độ rủi ro tín dụng. Trọng tâm của luận văn gồm so sánh hiệu năng mô hình, đánh giá ảnh hưởng của mất cân bằng lớp, phân tích vai trò của đặc trưng và đề xuất định hướng dữ liệu cho cảnh báo sớm.

Nghiên cứu sử dụng hai bộ dữ liệu tín dụng đã được ẩn danh từ một ngân hàng thương mại tại Việt Nam. Phạm vi dữ liệu bao gồm thông tin khách hàng, khoản vay và các yếu tố nghiệp vụ liên quan; tên ngân hàng và thông tin định danh trực tiếp không được công bố.

Các mô hình được đánh giá theo quy trình huấn luyện và kiểm định thống nhất trên từng bộ dữ liệu. Kết quả giữa hai bộ chỉ được đối chiếu về xu hướng do khác biệt về quy mô, cấu trúc biến và phân bố nhóm nợ. Các thông số dữ liệu, cấu hình tiền xử lý và thiết kế thực nghiệm được trình bày chi tiết tại Chương 3.

\section*{Động lực ứng dụng trí tuệ nhân tạo trong xếp hạng tín dụng}

Việc ứng dụng trí tuệ nhân tạo trong xếp hạng tín dụng được thúc đẩy bởi ba yêu cầu chính: khai thác tốt hơn dữ liệu ngày càng đa dạng; nâng cao tính nhất quán và hiệu quả của quy trình đánh giá; và hỗ trợ nhận diện sớm các khoản vay có dấu hiệu suy giảm. Trong lĩnh vực ngân hàng, hiệu quả dự báo cần đi cùng khả năng giải thích, kiểm soát và giám sát để mô hình có thể được sử dụng như một công cụ hỗ trợ quyết định.
